% GENERATED FILE. Edit canonical lessons, then run npm run generate:books.
% canonical-source-hash: fnv1a64:27f94e4de20cdf9a
% canonical-lessons: ZH-C538-miande, ZH-C538-conger, ZH-C538-benshen, ZH-C538-chansheng, ZH-C538-genben

\chapter{One Thing Leads to Another}
\label{ch:zh-one-thing-leads-to-another}
\hlchaptermodality{538}

\begin{chapteropening}
\textbf{By the end of this chapter:} \emph{I can say so as not to, and thereby, talk about the thing itself and what it gives rise to, and say not at all.}

Complete the last lesson of chapter 538: I can say so as not to, and thereby, talk about the thing itself and what it gives rise to, and say not at all.
\end{chapteropening}

\section[miǎnde]{miǎnde — so as not to; lest}

\label{lesson:ZH-C538-miande}



\begin{quote}
Before the new one: say the Mandarin for to cause damage, then the Mandarin for how well it works.
\end{quote}

\subsection*{You'll want to know: miǎnde}
\emph{miǎnde} — \textquotedblleft{}so as not to; lest\textquotedblright{}.

\begin{quote}
\emph{nǐ zuìhǎo dài yǔsǎn miǎnde línyǔ} — You'd better take an umbrella, so you don't get caught in the rain.
\end{quote}

Its characters come later; for now it is a word for the ear.

\begin{cousinweb}
\emph{miǎnde} is built from \emph{miǎn} 'to escape' and \zh{得}, both inherited from Middle Chinese, \emph{mjenX} and \emph{tok} (Guangyun, in Baxter's transcription), which glosses the first 'to stop' and 'to get free of'; the second is said light.
\end{cousinweb}

\begin{grammarlens}[title={What you've built}]
One more word for saying what would happen.
\end{grammarlens}

\subsection*{Your turn}
\begin{itemize}
\raggedright
  \item \emph{Say it:} \emph{miǎnde}
  \item \emph{Say it:} \emph{miǎnde}, once more
  \item \emph{Say it:} \emph{nǐ zuìhǎo dài yǔsǎn miǎnde línyǔ} — You'd better take an umbrella, so you don't get caught in the rain.
  \item \emph{Recall:} say the Mandarin for to cause damage, then the Mandarin for how well it works, then say \emph{miǎnde} again
\end{itemize}

\begin{tcolorbox}[breakable,colback=teal!4,colframe=teal!35!black,title={Before you move on}]
What does miǎnde mean? (\textquotedblleft{}So as not to; lest\textquotedblright{}.) Say it once more.
\end{tcolorbox}
\section[cóng'ér]{cóng'ér — and thereby; and in this way}

\label{lesson:ZH-C538-conger}



\begin{quote}
Before the new one: say the Mandarin for how well it works, then the Mandarin for so as not to.
\end{quote}

\subsection*{You'll want to know: cóng'ér}
\emph{cóng'ér} — \textquotedblleft{}and thereby; and in this way\textquotedblright{}.

\begin{quote}
\emph{wǒmen gǎi le shíjiān cóng'ér jiějué le zhège wèntí} — We changed the time, and so solved the problem.
\end{quote}

Its characters come later; for now it is a word for the ear.

\begin{cousinweb}
\emph{cóng'ér} is built from \zh{从} 'to follow' and \emph{ér} 'and', both inherited from Middle Chinese, \emph{dzjowng} and \emph{nyi} (Guangyun, in Baxter's transcription), which glosses the second as a grammar word.
\end{cousinweb}

\begin{grammarlens}[title={What you've built}]
One more word for saying what would happen.
\end{grammarlens}

\subsection*{Your turn}
\begin{itemize}
\raggedright
  \item \emph{Say it:} \emph{cóng'ér}
  \item \emph{Say it:} \emph{cóng'ér}, once more
  \item \emph{Say it:} \emph{wǒmen gǎi le shíjiān cóng'ér jiějué le zhège wèntí} — We changed the time, and so solved the problem.
  \item \emph{Recall:} say the Mandarin for how well it works, then the Mandarin for so as not to, then say \emph{cóng'ér} again
\end{itemize}

\begin{tcolorbox}[breakable,colback=teal!4,colframe=teal!35!black,title={Before you move on}]
What does cóng'ér mean? (\textquotedblleft{}And thereby; and in this way\textquotedblright{}.) Say it once more.
\end{tcolorbox}
\section[\texorpdfstring{\zh{本身} (běnshēn)}{本身 (běnshēn)}]{\zh{本身} (běnshēn) — itself (in itself)}

\label{lesson:ZH-C538-benshen}



\begin{quote}
Before the new one: say the Mandarin for so as not to, then the Mandarin for and thereby.
\end{quote}

\subsection*{You'll want to know: \zh{本身}}
\textbf{\zh{本身}} — \emph{běnshēn} — \textquotedblleft{}itself (in itself)\textquotedblright{}.

\begin{quote}
\textbf{\zh{这个问题本身不大}} \emph{zhège wèntí běnshēn bú dà} — The problem itself isn't a big one.
\end{quote}

\begin{cousinweb}
\textbf{\zh{本身}} (\emph{běnshēn}) is built from \zh{本} 'the root' and \zh{身} 'the body', both inherited from Middle Chinese, \emph{pwonX} and \emph{syin} (Guangyun, in Baxter's transcription), which says, after an older dictionary, that the base of a tree is \emph{pwonX}.
\end{cousinweb}

\begin{grammarlens}[title={What you've built}]
One more word for saying what would happen.
\end{grammarlens}

\subsection*{Your turn}
\begin{itemize}
\raggedright
  \item \emph{Say it:} \emph{běnshēn}
  \item \emph{Say it:} \emph{běnshēn}, once more
  \item \emph{Say it:} \emph{zhège wèntí běnshēn bú dà} — The problem itself isn't a big one.
  \item \emph{Recall:} say the Mandarin for so as not to, then the Mandarin for and thereby, then say \emph{běnshēn} again
\end{itemize}

\begin{tcolorbox}[breakable,colback=teal!4,colframe=teal!35!black,title={Before you move on}]
What does \zh{本身} mean? (\textquotedblleft{}Itself (in itself)\textquotedblright{}.) Say it once more.
\end{tcolorbox}
\section[chǎnshēng]{chǎnshēng — to arise; to come about}

\label{lesson:ZH-C538-chansheng}



\begin{quote}
Before the new one: say the Mandarin for and thereby, then the Mandarin for in itself.
\end{quote}

\subsection*{You'll want to know: chǎnshēng}
\emph{chǎnshēng} — \textquotedblleft{}to arise; to come about\textquotedblright{}.

\begin{quote}
\emph{yàoshi bù shuō qīngchu jiù huì chǎnshēng wùhuì} — If you don't make it clear, misunderstandings will arise.
\end{quote}

Its characters come later; for now it is a word for the ear.

\begin{cousinweb}
\emph{chǎnshēng} is built from \emph{chǎn} 'to bear' and \zh{生} 'to be born', both inherited from Middle Chinese, \emph{sreanX} and \emph{srjaeng} (Guangyun, in Baxter's transcription), which glosses the first 'to give birth'.
\end{cousinweb}

\begin{grammarlens}[title={What you've built}]
One more word for saying what would happen.
\end{grammarlens}

\subsection*{Your turn}
\begin{itemize}
\raggedright
  \item \emph{Say it:} \emph{chǎnshēng}
  \item \emph{Say it:} \emph{chǎnshēng}, once more
  \item \emph{Say it:} \emph{yàoshi bù shuō qīngchu jiù huì chǎnshēng wùhuì} — If you don't make it clear, misunderstandings will arise.
  \item \emph{Recall:} say the Mandarin for and thereby, then the Mandarin for in itself, then say \emph{chǎnshēng} again
\end{itemize}

\begin{tcolorbox}[breakable,colback=teal!4,colframe=teal!35!black,title={Before you move on}]
What does chǎnshēng mean? (\textquotedblleft{}To arise; to come about\textquotedblright{}.) Say it once more.
\end{tcolorbox}
\section[gēnběn]{gēnběn — not in the least; (not) at all}

\label{lesson:ZH-C538-genben}



\begin{quote}
Before the new one: say the Mandarin for in itself, then the Mandarin for to arise from something.
\end{quote}

\subsection*{You'll want to know: gēnběn}
\emph{gēnběn} — \textquotedblleft{}not in the least; (not) at all\textquotedblright{}.

\begin{quote}
\emph{jíshǐ nǐ shuō le tā yě gēnběn bù tīng} — Even if you tell him, he won't listen at all.
\end{quote}

Its characters come later; for now it is a word for the ear.

\begin{cousinweb}
\emph{gēnběn} is built from \emph{gēn} 'a root' and \zh{本} 'the base of a tree', both inherited from Middle Chinese, \emph{kon} and \emph{pwonX} (Guangyun, in Baxter's transcription), which glosses the first 'roots': down to the roots.
\end{cousinweb}

\begin{grammarlens}[title={What you've built}]
One more word for saying what would happen.
\end{grammarlens}

\subsection*{Your turn}
\begin{itemize}
\raggedright
  \item \emph{Say it:} \emph{gēnběn}
  \item \emph{Say it:} \emph{gēnběn}, once more
  \item \emph{Say it:} \emph{jíshǐ nǐ shuō le tā yě gēnběn bù tīng} — Even if you tell him, he won't listen at all.
  \item \emph{Recall:} say the Mandarin for in itself, then the Mandarin for to arise from something, then say \emph{gēnběn} again
\end{itemize}

\begin{tcolorbox}[breakable,colback=teal!4,colframe=teal!35!black,title={Before you move on}]
What does gēnběn mean? (\textquotedblleft{}Not in the least; (not) at all\textquotedblright{}.) Say it once more.
\end{tcolorbox}
